\section{An explicit geometric separator for the positive coefficient cone (NS-64)}
\label{sec:ns64-cone-separator}

The coefficient-cone exclusion in Proposition~\ref{prop:ns62-positive-cone}
has a small, explicit Hilbert-space witness. It quantifies that particular
failure without imposing any restriction on signed NB approximation.
Let $I=(1/6,\infty)$, put $p_1=\sigma_1|_I$ and
$p_n=\epsilon_n|_I$ for $2\leq n\leq6$, and write
\[
 G_I=(\langle p_m,p_n\rangle_{L^2(I)})_{m,n=1}^6,\quad
 \beta=(0,2,3/2,2/3,5/6,-1/5)^T,\quad
 d=G_I^{-1}e_6.
\]
The matrix $G_I$ is positive definite by the nested-support independence
proved in NS-62. In particular $d_6>0$.
The exact local target identity is $\tau|_I=\sum\beta_np_n$.

\begin{proposition}[A quantitative separation, confined to one coefficient cone]
\label{prop:ns64-cone-distance}
Define $\psi=\sum_{n=1}^6d_np_n$ on $I$ and extend it by zero outside
$I$. Then
\[
 \langle\psi,\sigma_1\rangle=0,\qquad
 \langle\psi,\epsilon_n\rangle=\mathbf1_{\{n=6\}}\quad(n\geq2),
 \qquad \langle\psi,\tau\rangle=-\frac15,
 \qquad \norm\psi^2=d_6.
\]
Consequently, if $\mathcal C_+$ denotes the closed cone generated by
$\sigma_1,\epsilon_2,\epsilon_3,\ldots$ with nonnegative real
coefficients, then
\begin{equation}
 \operatorname{dist}_H(\tau,\mathcal C_+)
 \geq\frac1{5\sqrt{d_6}}>0.00228323.
 \label{eq:ns64-positive-gap}
\end{equation}
The exact distance from $\tau|_I$ to the six-generator local cone is
$1/(5\sqrt{d_6})$; the interval replay encloses it between
$0.00228323766118601628$ and $0.00228323766118601629$.
The local equality is not asserted for the full-space distance.
\end{proposition}
\begin{proof}
The defining equation $G_Id=e_6$ gives the first six pairings and
$\norm\psi^2=d^TG_Id=d_6$. Every $\epsilon_n$, $n\geq7$, vanishes
on $I$. The target pairing follows from $\beta_6=-1/5$.
Thus $\langle\psi,f-\tau\rangle\geq1/5$ for every $f\in\mathcal C_+$;
continuity and Cauchy--Schwarz give the first inequality.

For the local equality, put
\[
 c=\beta+\frac{d}{5d_6},\qquad w=\frac{\psi}{5d_6}.
\]
Here $c_6=0$ exactly. The complete interval calculation below certifies
$c_1,\ldots,c_5>0$. Hence $\tau|_I+w=\sum c_np_n$ belongs to the
local cone. Its distance from $\tau|_I$ is $\norm w=1/(5\sqrt{d_6})$,
attaining the dual lower bound. Only the signs of these five finite
coefficients and the displayed decimal enclosure use certified
computation; the separating identity and formula for its bound are exact.
\end{proof}

\paragraph{Complete finite calculation, with no omitted tail.}
Under $t=1/x$, the domain $I$ becomes $0<t<6$ and its measure is
$dt/t^2$. Each $\sigma_n$ becomes $A(t/n)$. On the unit cell
$j<t<j+1$, $k=\lfloor j/n\rfloor$ and
\[
 A(t/n)=t/n-k\log t+k\log n+\log(k!).
\]
On $0<t<1$ the product $A(t/m)A(t/n)/t^2$ is exactly $1/(mn)$;
this cell includes the entire original $x>1$ tail. On the remaining
five cells the products have elementary primitives, retained in the
Maxima and Arb scripts. The triangular difference transformation gives
$G_I$ with every cross term present.

Arb at 256 and 384 bits, using certified preconditioned solves, gives
\[
 c\approx(0.00005096281,\ 1.99987424,\ 1.49456836,\
            0.68321323,\ 0.76905516,\ 0)^T,
\]
and
\[
 \frac1{25d_6}=0.000005213174217458189678\ldots.
\]
The figures in the vector are illustrative rounded displays of the
archived enclosures. All five positivity gates are strict, including
the small first coefficient. The two-precision comparison, complete
Gram and coefficient enclosures, and 39 exact Maxima checks are in
\texttt{evidence/v159/ns64/}. This is a separation of a specified
positive coefficient cone, not negative Weil energy, a lower bound for
the signed best-approximation error, or an obstruction to RH.
